%% Main manuscript for the double-word library paper.
\documentclass[sigplan,10pt,anonymous,review]{acmart}
\settopmatter{printfolios=true,printccs=false,printacmref=false}

\setcopyright{none}
\acmConference[CPP '27]{Certified Programs and Proofs}
  {January 11--12, 2027}{Mexico City, Mexico}

\bibliographystyle{ACM-Reference-Format}
\citestyle{acmauthoryear}

\usepackage{microtype}
\usepackage{amsmath,amsthm}
\usepackage{mathtools}
\usepackage{listings}
\usepackage{xcolor}
\usepackage{booktabs}
\usepackage{subcaption}
\usepackage{tikz}
\usetikzlibrary{arrows.meta,positioning,calc,fit}
\usepackage{hyperref}
\usepackage{cleveref}

\input{generated/evaluation_numbers}

\theoremstyle{plain}
\newtheorem{theorem}{Theorem}[section]
\newtheorem{lemma}[theorem]{Lemma}
\newtheorem{corollary}[theorem]{Corollary}
\newtheorem{proposition}[theorem]{Proposition}

\theoremstyle{definition}
\newtheorem{definition}[theorem]{Definition}
\newtheorem{example}[theorem]{Example}
\newtheorem{algorithm}[theorem]{Algorithm}
\crefname{definition}{Definition}{Definitions}
\crefname{example}{Example}{Examples}
\crefname{algorithm}{Algorithm}{Algorithms}

\theoremstyle{remark}
\newtheorem{remark}[theorem]{Remark}

\lstset{
  basicstyle=\ttfamily\small,
  keywordstyle=\bfseries\color{blue!70!black},
  commentstyle=\itshape\color{green!50!black},
  stringstyle=\color{orange!80!black},
  numbers=left,
  numberstyle=\tiny\color{gray},
  numbersep=6pt,
  breaklines=true,
  captionpos=b,
  frame=single,
  rulecolor=\color{gray!30},
  backgroundcolor=\color{gray!5},
}

\lstdefinelanguage{Coq}{
  keywords={Theorem,Lemma,Proof,Qed,Definition,Fixpoint,Inductive,
            forall,exists,fun,let,in,match,with,end,if,then,else,
            Variable,Hypothesis,Require,Import,Module,Section,
            apply,exact,intro,intros,rewrite,simpl,reflexivity,
            unfold,destruct,induction,case,split,left,right,
            specialize,assert,have,enough,lra,nra,ring,field},
  sensitive=true,
  morecomment=[l]{(*},
  morecomment=[s]{(*}{*)},
  morestring=[b]",
}

\newcommand{\fl}[1]{\operatorname{fl}(#1)}
\newcommand{\eps}{\varepsilon_{\mathrm{mach}}}
\newcommand{\RR}{\mathbb{R}}
\newcommand{\FF}{\mathbb{F}}
\newcommand{\ulp}{\operatorname{ulp}}
\newcommand{\TwoSum}{\textsc{TwoSum}}
\newcommand{\FastTwoSum}{\textsc{Fast2Sum}}
\newcommand{\FastTwoMult}{\textsc{Fast2Mult}}
\newcommand{\DWTimesFP}{\textsc{DWTimesFP}}
\newcommand{\DWDivFP}{\textsc{DWDivFP}}
\newcommand{\DWTimesFPCond}{\mathsf{Cond}_{\mathrm{DWTimesFP}}}
\newcommand{\DWArithOp}{\textsc{DWArithOp}}
\newcommand{\DWArithOpCond}{\mathsf{Cond}_{\mathrm{DWArithOp}}}
\newcommand{\DW}{\mathrm{DW}}

\author{Ariel Kellison}
\affiliation{%
  \institution{Code Metal, Inc.}
  \city{Boston}
  \state{Massachusetts}
  \country{USA}
}

\title{A Verified Library of Double-Double Arithmetic}

\begin{abstract}
Double-double arithmetic represents a real number as the unevaluated sum of two
binary64 values, providing roughly twice the precision of binary64.  Existing
analyses prove error and validity guarantees for its algorithms only under
explicit conditions that conventional libraries do not enforce.

We present a verified C library of double-double arithmetic consisting of
three error-free transformations and six double-word operations.  Our Rocq
development extends prior analysis that assumes an unbounded exponent range
to IEEE~754 arithmetic, accounting for overflow and gradual underflow.  Using
VST, we verify that the C kernels implement the analyzed algorithms and that
runtime checks establish the input conditions and result finiteness required
for the proved guarantees.  Thus, on normal return, the public double-word
functions return valid double-doubles within their proved error bounds, while
the error-free transformations guarantee exact reconstruction.

Across tests of the six double-word operations spanning the binary64 exponent
range, the result of every accepted call was valid and within its proved error
bound.  The largest observed error was \RQOneMaximumBoundRatio{} times the
corresponding bound.  To assess the usefulness of the runtime checks, we
compare with QD, an established package for double-double arithmetic.  For
calls rejected by the six default double-word functions, QD returned a finite,
valid result within the corresponding library bound in only
\RQTwoDefaultQDWithinPercent{} of cases.  For rejected calls to the
\FastTwoSum{} and \FastTwoMult{} error-free transformations, it reconstructed
the exact result in only \RQTwoFastSumExactPercent{} and
\RQTwoFastMultExactPercent{} of cases.  With CompCert, the runtime checks added
\RQThreeMinimumAddedTime{}--\RQThreeMaximumAddedTime{} nanoseconds per call
relative to the unchecked kernels.  With Clang, the public functions were
\RQThreeMinimumQDRatio{}--\RQThreeMaximumQDRatio{} times slower than QD.
\end{abstract}

\begin{CCSXML}
<ccs2012>
 <concept>
  <concept_id>10003752.10003790.10003794</concept_id>
  <concept_desc>Theory of computation~Automated reasoning</concept_desc>
  <concept_significance>500</concept_significance>
 </concept>
 <concept>
  <concept_id>10010583.10010786.10010813</concept_id>
  <concept_desc>Hardware~Arithmetic and datapath circuits</concept_desc>
  <concept_significance>300</concept_significance>
 </concept>
</ccs2012>
\end{CCSXML}

\ccsdesc[500]{Theory of computation~Automated reasoning}
\ccsdesc[300]{Hardware~Arithmetic and datapath circuits}

\keywords{floating-point arithmetic, double-word arithmetic, formal verification,
  Rocq, error analysis, IEEE 754}

\begin{document}
\maketitle

\section{Introduction}
\label{sec:intro}

Binary64 arithmetic is fast and widely available, but its precision is not
enough for every numerical computation \citep{Higham2002}.  Arbitrary-precision
libraries such as MPFR address this limitation by allowing each variable to
have an independently chosen precision \citep{Fousse2007}.  When roughly twice
the precision of binary64 is sufficient, double-double arithmetic offers a
more specialized alternative.  It represents a number as the unevaluated sum
of two binary64 values and performs arithmetic on the pair.  This approach has
been used to improve reproducibility in parallel climate models
and to provide additional precision in quantum electrodynamic calculations
\citep{HeDing2001,Aoyama2015}.

The apparent simplicity of the representation hides an important hazard.
\citet{Li2002} described double-double arithmetic as ``an `attractive
nuisance,' like an unfenced backyard swimming pool.''  A pair of binary64
values does not form an ordinary wider IEEE~754 format, and a double-double
operation is not obtained by computing an exact result and rounding it once.
Instead, a double-double operation is implemented as a sequence of rounded
binary64 computations.  Different algorithms can implement the same
double-double operation, and the error bound depends on the chosen algorithm.
A separate argument must also show that the returned pair is a valid
double-double.  The error bound and validity guarantee hold only under
conditions on the operands and intermediate results.

We present a verified C library whose public C interface enforces these
conditions.  The library provides three error-free transformations and six
double-double operations.  On normal return, each double-double operation produces 
a valid double-double within its proved error bound, while each error-free 
transformation guarantees exact reconstruction.

\subsection{Double-Word Arithmetic and Its Numerical Guarantees}
\label{sec:intro-double-word}

We state the numerical analysis for double-word arithmetic over a generic
underlying floating-point format.  Double-double arithmetic is the binary64
instance of double-word arithmetic.  A double word is a floating-point
expansion with two components \citep{Dekker1971,Priest1991}.  Let
$\mathbb{F}$ denote the underlying format, and let $\operatorname{RN}$ denote
round-to-nearest, ties-to-even. Following
\citet[Definition~1.4]{Joldes2017}, we use the following definition.

\begin{definition}[Double word]
\label[definition]{def:double-word}
A double word $x$ over $\mathbb{F}$ is the unevaluated sum
\begin{equation}
  x=x_h+x_\ell
\end{equation}
of two values $x_h,x_\ell\in\mathbb{F}$ such that
\begin{equation}
  x_h=\operatorname{RN}(x).
  \label{eq:dw-representation}
\end{equation}
\end{definition}

Thus, rounding the represented value $x$ recovers the high component $x_h$,
while the low component is the exact residual $x_\ell=x-x_h$.  For an exact
real result $r$, a double-word algorithm should return $(z_h,z_\ell)$
satisfying two distinct properties.  Writing $\gamma$ for an error bound
derived by the analysis, the accuracy property is
\begin{equation}
  \left|(z_h+z_\ell)-r\right|\leq\gamma,
  \label{eq:intro-accuracy-property}
\end{equation}
and the validity property is that the returned components form a double word
$z$ according to \cref{def:double-word}.  The validity property supplies the
representation precondition required when the result is used by a later
double-word operation.

The numerical results on which we build
\citep{Joldes2017,Rideau2022} establish both properties for double-word
addition, multiplication, and division when underflow and overflow do not
occur.  \citet{Rideau2022} mechanize these results using Flocq's FLX format
\citep{Boldo2011}, whose unbounded exponent range represents this assumption.

Native IEEE~754 arithmetic requires both ends of the exponent range to be
considered.  Near the upper end, an intermediate computation can overflow even when its
operands are finite, and the returned pair may consequently contain an
infinity or NaN.  Neither an infinity nor a NaN represents a real value, so
such a pair can satisfy neither the accuracy property nor the valid
double-word property.  Both properties
therefore require the returned components to be finite.  Near the lower end,
the fixed spacing of subnormal values can invalidate exact relations used by
the basic building blocks of double-word algorithms and prevent rounding
errors from being bounded solely relative to the exact result.

Our analysis determines how these effects change the guarantee for each
operation.  Addition retains the relative bounds of the earlier analysis.
Multiplication and division instead satisfy mixed bounds: a relative term
accounts for ordinary rounding error, while an absolute term accounts for
errors introduced near the smallest representable values.  For every
operation, the resulting theorem gives sufficient conditions under which the
returned pair is a valid double word and its represented value satisfies the
stated error bound.

Conditions governing internal computations must be established before
execution.  Result finiteness can be established in either of two
ways: stronger conditions on the inputs can imply it before execution, or the
returned components can be checked afterward.  Making the numerical results
available through ordinary C calls therefore requires runtime checks that
establish the appropriate conditions before or after the kernel call.

\subsection{This Work: A Verified C Library for Double-Double Arithmetic}

For each of the six double-double operations, the library provides two public
wrappers around the same arithmetic kernel.  The default wrapper performs a
weak-precondition check before execution, invokes the kernel, and then checks
that both returned components are finite.  The other wrapper performs a
strong-precondition check that establishes result finiteness before execution,
so no result check is needed afterward.  The terms \emph{weak} and
\emph{strong} describe the precondition checks, not the wrappers.  Both
wrappers provide the same contract on normal return: the represented result
satisfies the proved error bound, and the returned pair is a valid double word
over binary64---that is, a valid double-double result.

For each operation, we define in Rocq a functional model over a generic
radix-two format and prove its numerical guarantee using Flocq.  We then use
VST to prove that the C kernel implements the binary64 instance of this model
and that the runtime checks establish the conditions required by its
numerical theorem.  Together, these proofs establish the contracts of the
public C interface.

We evaluate whether the resulting checks are both effective and practical.
Across tests spanning the binary64 exponent range, every double-double result
returned normally was valid and within its proved error bound; the largest
observed error was \RQOneMaximumBoundRatio{} times the corresponding bound.
For calls stopped by a runtime check, the experiment ran the corresponding QD
operation on the same inputs and measured whether its result satisfied the
corresponding numerical criterion.  Finally,
the performance measurements separate the cost of the runtime checks from the
cost of the arithmetic kernels.  \Cref{sec:evaluation} gives the experimental
methods and results.

\subsection{Contributions}

In summary, this paper makes four contributions:

\begin{itemize}
  \item \textbf{Runtime checks for double-double arithmetic
        (\Cref{sec:overview}).}
        We implement public wrappers whose runtime checks establish the
        sufficient conditions required by these numerical results.

  \item \textbf{Numerical guarantees over IEEE~754
        (\Cref{sec:numerics}).}
        We extend prior mechanized analysis to the finite exponent range of
        IEEE~754, proving error bounds and valid double-word results for six
        operations and guarantees of exact reconstruction for three
        error-free transformations.

  \item \textbf{A verified C implementation of double-double arithmetic
        (\Cref{sec:mechanization}).}
        Using VST, we connect the Rocq functional models and numerical
        theorems to the binary64 C kernels, runtime checks, and contracts of
        the public C interface.

  \item \textbf{Empirical evaluation
        (\Cref{sec:evaluation}).}
        We evaluate the tightness of the proved bounds, the behavior of QD on
        inputs associated with calls for which a runtime check failed, and the
        cost of the public wrappers relative to the arithmetic kernels and QD.
\end{itemize}

The accompanying artifact contains the C sources, Rocq and VST proofs, and
evaluation scripts at \url{ARTIFACT-URL}, revision
\texttt{ARTIFACT-COMMIT}.

\section{Runtime Checks for Double-Double Arithmetic}
\label{sec:overview}

\lstdefinestyle{sectiontwoc}{
  language=C,
  basicstyle=\ttfamily\normalsize,
  columns=fullflexible,
  keepspaces=true,
  numbers=left,
  numberstyle=\scriptsize\color{gray},
  numbersep=6pt,
  breaklines=true,
  breakatwhitespace=true,
  frame=single,
  framesep=3pt,
  xleftmargin=1.5em,
  framexleftmargin=1.2em,
  framexrightmargin=3pt,
  aboveskip=0pt,
  belowskip=0pt,
  showstringspaces=false
}

This section uses division of a double-double number by a binary64 value to
illustrate a library interface.  We present its two public wrappers, relate
their C data type to the mathematical definition of a double word, and explain
their shared kernel and runtime checks.  We then state their common contract
and summarize the remaining library operations.

\subsection{Runtime Checks for Dividing a Double-Double by a Binary64 Value}
\label{sec:overview-division-wrappers}

The C library provides two public wrappers for dividing a double-double
number by a binary64 value.  Both call the same arithmetic kernel and use the
\texttt{struct dword} type shown in \Cref{fig:dword-struct} to store their
double-double operands and results.

\begin{figure}[t]
  \begin{lstlisting}[style=sectiontwoc]
struct dword {
    double s;
    double t;
};
\end{lstlisting}
  \caption{The C structure used to store a double-double number.  The fields
  contain its two binary64 components.}
  \Description{A C structure containing two binary64 fields named s and t.}
  \label{fig:dword-struct}
\end{figure}

The default public function is the weak-precondition wrapper shown in
\Cref{fig:weak-precondition-division-wrapper}.  The argument \texttt{x}
points to the stored dividend, \texttt{y} is the binary64 divisor, and
\texttt{z} points to storage for the quotient.  Before invoking the kernel,
the wrapper calls
\texttt{dw\_div\_fp\_\allowbreak weak\_precond\_\allowbreak check}.  After the
kernel writes the quotient to \texttt{z}, it calls
\texttt{isfinite\_dword} to check the computed pair.  If either runtime check
fails, the function exits.

\begin{figure}[t]
  \begin{lstlisting}[style=sectiontwoc]
void dw_div_fp(
    struct dword *z, struct dword *x, double y)
{
    if (!dw_div_fp_weak_precond_check(x, y))
        exit(EXIT_FAILURE);

    dw_div_fp_kernel(z, x, y);

    if (!isfinite_dword(z))
        exit(EXIT_FAILURE);
}
\end{lstlisting}
  \caption{The weak-precondition division wrapper.  It checks the operands
  before invoking the kernel and checks result finiteness afterward.}
  \Description{A C function that checks its operands, computes a
  double-double quotient, and checks both result components for finiteness.}
  \label{fig:weak-precondition-division-wrapper}
\end{figure}

Some callers may want the runtime checks to establish result finiteness before
the arithmetic kernel is executed.  The strong-precondition wrapper in
\Cref{fig:strong-precondition-division-wrapper} calls
\texttt{dw\_div\_fp\_\allowbreak strong\_precond\_\allowbreak check} for this
purpose.  If the check
fails, the function exits without invoking the kernel.  If it succeeds, the
kernel writes the quotient directly to \texttt{z}; no result check is needed
afterward.

\begin{figure}[t]
  \begin{lstlisting}[style=sectiontwoc]
void dw_div_fp_strong(
    struct dword *z, struct dword *x, double y)
{
    if (!dw_div_fp_strong_precond_check(x, y))
        exit(EXIT_FAILURE);

    dw_div_fp_kernel(z, x, y);
}
\end{lstlisting}
  \caption{The strong-precondition division wrapper.  Its runtime check
  establishes the required conditions, including result finiteness, before
  the kernel is invoked.}
  \Description{A C function that exits when its strong precondition check
  fails and otherwise calls the division kernel with the caller-provided
  output structure.}
  \label{fig:strong-precondition-division-wrapper}
\end{figure}

The following subsections explain the stored representation, the shared
kernel, and the runtime checks used by these two functions.

\subsection{The Double-Double Data Type}
\label{sec:overview-values}

% Connect struct dword, which is specific to the binary64 C library, to the
% general definition of a double word in Section 1.  Keep the distinction
% explicit: the C fields store the binary64 components; the cited mathematical
% relation determines whether their unevaluated sum is a valid double word.
% End by motivating the C kernel that computes the returned pair.  Do not
% promise a displayed generic algorithm or functional model in Section 2.

\Cref{def:double-word} defines a double word over a generic underlying
floating-point format $\mathbb{F}$.  The C library instantiates $\mathbb{F}$
as binary64 and uses round-to-nearest, ties-to-even.  A double word over
binary64 is a double-double number.  The fields \texttt{s} and \texttt{t} in
\Cref{fig:dword-struct} store its high and low words, respectively.

The C type does not itself enforce the double-word relation.  Because the
numerical guarantee assumes that the dividend is a valid double word, both
wrappers check this property before calling the arithmetic kernel.
\Cref{ex:struct-dword} illustrates how the stored components represent a value
and how this relation is checked.

\begin{example}[Checking a \texttt{struct dword}]
\label[example]{ex:struct-dword}
The following object stores the binary64 values $1$ and $2^{-54}$:

\begin{lstlisting}[language=C]
struct dword x = {
    .s = 1.0,
    .t = 0x1p-54
};
\end{lstlisting}

The value represented by this object is the unevaluated real sum
$1+2^{-54}$.  Under binary64 round-to-nearest, ties-to-even,
$\operatorname{RN}(1+2^{-54})=1$, so the stored pair is a valid double word according to
\Cref{def:double-word}.  Its represented value is not merely the value stored
in \texttt{x.s}; the field \texttt{x.t} retains the correction $2^{-54}$.

The runtime checks used by the wrappers first require both fields
to be finite and then evaluate
\begin{lstlisting}[language=C,numbers=none]
x->s + x->t == x->s
\end{lstlisting}
using binary64 arithmetic.  This comparison succeeds for the object above
because the computed sum rounds to \texttt{x.s}.  Soundness of the check
shows that a successful comparison, together with the finiteness tests,
establishes the mathematical relation in
\Cref{eq:dw-representation}.  In contrast, replacing \texttt{0x1p-54} with
\texttt{0x1p-52} makes the computed sum different from \texttt{x.s}, so the
check fails for the pair.
\end{example}

This test establishes one of the input preconditions required for the
numerical guarantee of the division kernel.  Division has additional
preconditions involving the range of the operands and computations performed
within the kernel.  Their purpose is best understood from the implementation,
so we present the shared C kernel before examining the complete checks
performed by the two wrappers.

\subsection{The Division Kernel}
\label{sec:overview-division-kernel}

The weak-precondition wrapper in
\Cref{fig:weak-precondition-division-wrapper} and the strong-precondition
wrapper in \Cref{fig:strong-precondition-division-wrapper} call the same
arithmetic kernel.
\Cref{fig:dw-div-fp-kernel} shows \texttt{dw\_div\_fp\_kernel}, the C
implementation of the division method given as Algorithm~15 by
\citet{Joldes2017}.

\begin{figure}[t]
  \begin{lstlisting}[style=sectiontwoc]
void dw_div_fp_kernel(
    struct dword *z, struct dword *x, double y)
{
    double th;
    struct dword p;
    double dh;
    double dt;
    double d;
    double tl;

    th = x->s / y;
    fast_2mult(&p, th, y);
    dh = x->s - p.s;
    dt = dh - p.t;
    d = dt + x->t;
    tl = d / y;
    fast_2sum(z, th, tl);
}
\end{lstlisting}
  \caption{The binary64 C kernel for dividing a double-double number by a
  binary64 value.}
  \Description{A C function that computes a leading quotient, forms a
  correction from the reconstructed product and the low word of the dividend,
  and combines the leading quotient and correction using FastTwoSum.}
  \label{fig:dw-div-fp-kernel}
\end{figure}

The kernel first computes the leading quotient \texttt{th} from the high word
of the dividend.  The call to \texttt{fast\_2mult} multiplies this quotient by
the divisor and returns two components used to reconstruct the product.  The
computations of \texttt{dh}, \texttt{dt}, and \texttt{d} use this product and
the low word of the dividend to form a residual.  Dividing the residual by
\texttt{y} produces the quotient correction \texttt{tl}.  Finally,
\texttt{fast\_2sum} combines the leading quotient and correction into the pair
stored in \texttt{z}.

\Cref{sec:numerics} presents \texttt{DWDivFP}, the generic Rocq functional
model that specifies this computation, and proves its numerical guarantees
using Flocq.  As described in \Cref{sec:kernel-refinement}, a VST refinement
proof establishes that \texttt{dw\_div\_fp\_kernel} implements the binary64
instance of \texttt{DWDivFP}.

The refinement result identifies the pair computed by C, but the accuracy and
valid-double-word guarantees require additional mathematical conditions.  The
two wrappers establish these conditions differently.  The following
subsections present their executable C checks, beginning with the
weak-precondition wrapper in
\Cref{fig:weak-precondition-division-wrapper}, which checks the computed
result for finiteness.

\subsection{Checking Finiteness after Division}
\label{sec:overview-post-check}

The weak-precondition wrapper uses the helper functions in
\Cref{fig:finite-checks} to test the operands and the pair computed by the
kernel.  Under strict IEEE~754 semantics, $x-x$ is zero when $x$ is finite
and NaN when $x$ is an infinity or NaN.  Because comparison with NaN returns
false, \texttt{isfinite\_double} returns $1$ exactly when its argument is
finite.  The function \texttt{isfinite\_dword} applies this test to both
stored components.  The complete input check is shown in
\Cref{fig:dw-div-fp-input-check}.

\begin{figure}[t]
  \begin{lstlisting}[style=sectiontwoc]
int isfinite_double(double x)
{
    return (x - x) == 0.0;
}

int isfinite_dword(
    const struct dword *x)
{
    return isfinite_double(x->s) &&
           isfinite_double(x->t);
}
\end{lstlisting}
  \caption{The C functions used to check a binary64 value and both components
  of a \texttt{struct dword} for finiteness.}
  \Description{Two C functions.  The first tests whether a binary64 value is
  finite by subtracting it from itself and comparing the result with zero.
  The second applies this test to both fields of a double-double value.}
  \label{fig:finite-checks}
\end{figure}

\begin{figure}[t]
  \begin{lstlisting}[style=sectiontwoc]
int dw_div_fp_weak_precond_check(
    struct dword *x, double y)
{
    int fin_xh = isfinite_double(x->s);
    int fin_xl = isfinite_double(x->t);
    int fin_y  = isfinite_double(y);

    int dwx = fin_xh && fin_xl &&
              (x->s + x->t == x->s);

    double q_plo = ldexp(1.0, -968);
    int xh_lo = (x->s == 0.0) ||
                (q_plo <= fabs(x->s));

    return dwx && fin_y && xh_lo;
}
\end{lstlisting}
  \caption{The input check used by the weak-precondition division wrapper.
  It tests operand finiteness, dividend validity, and the lower bound required
  by the kernel.}
  \Description{A C function that checks the finiteness of the operands, the
  double-word relation for the dividend, and either a zero high word or a
  lower bound on its magnitude.}
  \label{fig:dw-div-fp-input-check}
\end{figure}

The clause \texttt{dwx} combines finiteness tests for the two stored
components with the comparison illustrated in \Cref{ex:struct-dword}.
The check also tests the divisor for finiteness.  These finiteness tests provide
early rejection for infinities and NaNs; they are not separate hypotheses of
the numerical theorem.

The remaining condition permits $x_h=0$ or requires
$|x_h|\geq2^{-968}$.  In the first case, the double-word relation
checked by \texttt{dwx} also implies $x_\ell=0$, so the dividend is exactly
zero.  For a nonzero high word, the lower bound is required by the product
reconstruction performed by \texttt{fast\_2mult} in
\Cref{fig:dw-div-fp-kernel}.  \Cref{sec:division-numerics} derives this
requirement from the numerical analysis.

The input check imposes no upper bound on the operands or leading quotient.
After the kernel executes, \texttt{isfinite\_\allowbreak dword} checks both
components of the computed pair.  It succeeds for normal values, subnormal
values, and zeros, and fails if either component is an infinity or NaN.

\subsection{Guaranteeing Finiteness before Division}
\label{sec:overview-strong-precondition-check}

The strong-precondition wrapper extends the preceding input check with
bounds that establish result finiteness before the kernel executes.  Its
complete check is shown in \Cref{fig:dw-div-fp-check}.  The additional
comparisons bound the high word of the dividend and the rounded leading
quotient.

\begin{figure}[t]
  \begin{lstlisting}[style=sectiontwoc]
int dw_div_fp_strong_precond_check(
    struct dword *x, double y)
{
    int fin_xh = isfinite_double(x->s);
    int fin_xl = isfinite_double(x->t);
    int fin_y  = isfinite_double(y);

    int dwx = fin_xh && fin_xl &&
              (x->s + x->t == x->s);

    double q_omega = ldexp(1.0, 1022);
    double q_plo   = ldexp(1.0, -968);

    int xh_hi = fabs(x->s) <= q_omega;
    int xh_lo = (x->s == 0.0) ||
                (q_plo <= fabs(x->s));

    double q = x->s / y;
    int qt_hi = fabs(q) < q_omega;

    return dwx && fin_y &&
           xh_hi && xh_lo && qt_hi;
}
\end{lstlisting}
  \caption{The input check used by the strong-precondition division wrapper.
  In addition to the tests in \Cref{fig:dw-div-fp-input-check}, it bounds the
  high word and rounded leading quotient.}
  \Description{A C function that checks the operands and requires upper
  bounds on the high word and rounded leading quotient before allowing the
  division kernel to execute.}
  \label{fig:dw-div-fp-check}
\end{figure}

The first additional test requires $|x_h|\leq 2^{1022}$.  The second computes
the same leading quotient as the first line of
\texttt{dw\_div\_fp\_kernel}, $t_h=\operatorname{RN}(x_h/y)$, and requires
$|t_h|<2^{1022}$.  The C variable \texttt{q} stores this rounded
leading quotient.  In particular, the strict comparison fails for a nonfinite
leading quotient.  The two upper bounds are sufficient to show before
execution that the computations in the kernel and both components of its
result remain finite.  No lower bound is imposed on the leading quotient, so
the check permits subnormal values.

These bounds are sufficient but not necessary.  Checks based on them may
therefore fail for operands whose kernel result would be a finite pair.  This
conservative rejection is the cost of establishing finiteness before execution.
\Cref{sec:numerics} derives the bounds and connects them to the generic
numerical conditions.

\subsection{The Division Contract}
\label{sec:overview-division-guarantee}

Although the weak-precondition wrapper in
\Cref{fig:weak-precondition-division-wrapper} and the strong-precondition
wrapper in \Cref{fig:strong-precondition-division-wrapper} establish
result finiteness at different points, they provide the same contract on
normal return.

\begin{theorem}[Normal-Return Guarantee for the Division Wrappers]
\label{thm:guarded-division}
Suppose that \texttt{x} initially stores $(x_h,x_\ell)$ and that a call to
either \texttt{dw\_div\_fp} or
\texttt{dw\_div\_fp\_\allowbreak strong} returns normally
with \texttt{z} storing $(z_h,z_\ell)$.  Writing
$q=(x_h+x_\ell)/y$ for the exact real quotient, the represented result
satisfies
\[
  \left|(z_h+z_\ell)-q\right|
  \leq 3u^2|q|+\eta,
\]
where $u=2^{-53}$ is the unit roundoff and $\eta=2^{-1074}$ is the smallest
positive subnormal value for binary64.  The returned pair is also a valid
double-double according to \Cref{def:double-word}.
\end{theorem}

Callers need not supply numerical preconditions.  Instead, each wrapper
establishes the sufficient conditions required by the underlying
numerical result.  The guarantee is conditional on normal return: if any
check fails, the wrapper follows its nonreturning error path and provides no
numerical result.

\Cref{sec:mechanization} shows how the VST proofs combine soundness of the
corresponding check, refinement of the C kernel, and the numerical theorem to
establish this contract.

\subsection{Operations and Guarantees of the C Library}
\label{sec:overview-library-guarantees}

The two division wrappers illustrate the form of the contract provided for
each double-word operation in the C library.  \Cref{tab:operation-preview}
places division alongside the other available operations and summarizes the
accuracy guarantee provided when a wrapper returns normally.

\begin{table*}[t]
\centering
\caption{Normal-return accuracy guarantees provided by the C library.  Here
$a$ and $b$ are binary64 operands, $x$ and $y$ are the values represented by
double-double operands, $r$ is the exact result, and
$\widehat r=z_h+z_\ell$ is the value represented by the returned pair.  For
binary64, the unit roundoff is $u=2^{-53}$ and the smallest positive
subnormal value is $\eta=2^{-1074}$.  Every returned pair is also a valid
double-double number.}
\Description{The operations provided by the C library, their exact real
results, and their accuracy guarantees when a wrapper returns normally.}
\label{tab:operation-preview}
\small
\begin{tabular}{@{}p{0.49\textwidth}p{0.45\textwidth}@{}}
\toprule
Operation & Normal-return accuracy guarantee \\
\midrule
\multicolumn{2}{@{}l}{\emph{Error-free transformations}} \\[2pt]
\TwoSum{}$(a,b)$ and \FastTwoSum{}$(a,b)$; $r=a+b$
  & $\widehat r=r$ \\
\FastTwoMult{}$(a,b)$; $r=ab$
  & $\widehat r=r$ \\
\addlinespace
\multicolumn{2}{@{}l}{\emph{Double-word operations}} \\[2pt]
Addition $(x,a)$; $r=x+a$
  & $|\widehat r-r|\leq 2u^2|r|$ \\
Addition $(x,y)$; $r=x+y$
  & $|\widehat r-r|
       \leq \dfrac{3u^2}{1-4u}|r|$ \\
Multiplication $(x,a)$; $r=xa$
  & $|\widehat r-r|
       \leq 2u^2|r|+2.51\eta$ \\
Multiplication $(x,y)$; $r=xy$
  & $|\widehat r-r|
       \leq 4u^2|r|+2.5\eta$ \\
Division $(x,a)$; $r=x/a$
  & $|\widehat r-r|
       \leq 3u^2|r|+\eta$ \\
Division $(x,y)$; $r=x/y$
  & $|\widehat r-r|
       \leq 9.8u^2|r|+2.5\eta$ \\
\bottomrule
\end{tabular}
\end{table*}

Each of the six double-word operations in \Cref{tab:operation-preview} is
available through both wrappers.  The two wrappers provide the same
guarantee shown in the table, although their checks can succeed on different
sets of inputs.
Each error-free transformation instead has one wrapper for its
exactness guarantee.

The set of inputs for which all runtime checks succeed depends on which
computations in the corresponding algorithm must remain finite and
how rounding near the lower end of the format affects its result.  Determining
these domains therefore requires a numerical analysis that accounts for the
finite exponent range of IEEE~754 arithmetic.  \Cref{sec:numerics} develops
this analysis and derives sufficient conditions for each error bound and valid
double-word result.

\section{Numerical Guarantees for Double-Word Arithmetic over IEEE~754}
\label{sec:numerics}

To prove the numerical guarantees presented in \Cref{sec:overview},
we define \emph{functional models} of the arithmetic algorithms.
These mathematical models are expressed as Rocq functions using
IEEE~754 operations from VCFloat2 \citep{AppelKellison2024} and the
binary floating-point representation from Flocq \citep{Boldo2011}.
Each model follows the sequence of rounded operations in the
corresponding arithmetic kernel and returns a pair of floating-point
values, independently of the C representation. We define the models
over a generic radix-two format; \FenceDD{} uses the binary64
instances.

We first establish the rounding properties and error-free
transformation results needed to analyze complete double-word
operations. We then continue the division example from
\Cref{sec:overview-division}, giving a proof sketch of the numerical
guarantee and explaining how the proofs for the remaining operations
differ.

\subsection{FLX, FLT, and IEEE~754 Arithmetic}
\label{sec:flx-to-ieee}

To extend the earlier mechanization by \citeauthor{Rideau2022} to
IEEE~754 arithmetic \citep{IEEE754-2019}, we relate rounding in two
Flocq formats with the same radix and precision: floating-point with
unbounded exponents (FLX) and floating-point with gradual underflow
(FLT). The earlier mechanization uses FLX and therefore excludes
both underflow and overflow. FLT imposes a lower exponent bound and
includes subnormal values.

For radix $\beta>1$, precision $p>0$, and lower exponent $e_{\min}$, the two
formats represent
\[
\begin{aligned}
  \mathbb{F}_{\mathrm{FLX}}(\beta,p)
    &=
    \left\{m\beta^q \mathrel{}\middle|\mathrel{}
      m,q\in\mathbb{Z},\ |m|<\beta^p\right\},\\
  \mathbb{F}_{\mathrm{FLT}}(\beta,p,e_{\min})
    &=
    \left\{m\beta^q \mathrel{}\middle|\mathrel{}
      m,q\in\mathbb{Z},\ |m|<\beta^p,\ q\geq e_{\min}\right\}.
\end{aligned}
\]
Thus,
$\mathbb{F}_{\mathrm{FLT}}(\beta,p,e_{\min})
 \subseteq\mathbb{F}_{\mathrm{FLX}}(\beta,p)$.

The lower exponent bound determines the smallest positive
representable value $\eta$ and the smallest positive normal value $\nu$:
\begin{equation}
  \eta=\beta^{e_{\min}}
  \qquad\text{and}\qquad
  \nu=\beta^{e_{\min}+p-1}.
  \label{eq:flt-thresholds}
\end{equation}
Nonzero representable values with magnitude below $\nu$ are
subnormal. In this range, adjacent representable values are spaced
$\eta$ apart. For binary64, $\eta=2^{-1074}$ and $\nu=2^{-1022}$.

Let $\operatorname{RN}_{\mathrm{FLX}}$ and
$\operatorname{RN}_{\mathrm{FLT}}$ denote rounding to nearest with
ties to even in the two formats. For exact real values with
magnitude at least $\nu$, the Flocq theorem
\texttt{FLT.round\_FLT\_FLX} gives
\begin{equation}
  |x|\geq\nu
  \quad\Longrightarrow\quad
  \operatorname{RN}_{\mathrm{FLT}}(x)
    =\operatorname{RN}_{\mathrm{FLX}}(x).
  \label{eq:round-flt-flx}
\end{equation}
This equality allows us to reuse the earlier FLX arguments for
operations whose exact results have magnitude at least $\nu$.
For smaller exact results, the two roundings may differ, so we
analyze FLT rounding directly.

Neither FLX nor FLT has an upper exponent bound. Our error analysis
assumes that the returned components are finite and uses this
hypothesis to relate the IEEE~754 computations to FLT rounding.
We separately establish sufficient conditions under which the
returned components are finite.

\subsection{Error-Free Transformations over IEEE~754}
\label{sec:error-free-transformations}

Double-word algorithms use error-free transformations that return
a rounded sum or product together with a correction intended to
equal the exact rounding error. To extend the analysis of these
algorithms to IEEE~754 arithmetic, we must determine whether this
correction remains exact when gradual underflow occurs. For
addition, the equality remains valid under suitable conditions.
For multiplication, the exact rounding error may be too small to
represent. We identify a sufficient condition under which the
correction remains exact and otherwise bound the error in the sum
of the returned components.

\begin{definition}[Error-free transformation]
\label[definition]{def:error-free-transformation}
For an arithmetic operation with exact real result $r$, an error-free
transformation returns a pair $(r_h,r_\ell)$ satisfying
\[
  r_h=\operatorname{RN}(r)
  \qquad\text{and}\qquad
  r=r_h+r_\ell.
\]
\end{definition}

Thus, $r_\ell=r-r_h$ is the exact rounding error. Every pair returned by an
error-free transformation also forms a double word because
$r_h=\operatorname{RN}(r_h+r_\ell)$ follows from
\Cref{def:error-free-transformation}.

\Cref{fig:error-free-transformations} gives the three error-free
transformations used by the double-word algorithms in \FenceDD{}. 
We use the naming conventions and algorithms from 
\citet[Algorithms\ 1--3]{Joldes2017}.
\TwoSum{} is due to Knuth, while
\FastTwoSum{} and the original error-free transformations for products are due
to Dekker \citep{Knuth1997,Dekker1971}.

\begin{figure}[t]
  \centering
  \begingroup
  \setlength{\fboxsep}{6pt}
  \setlength{\fboxrule}{0.4pt}
  \begin{subfigure}[t]{\columnwidth}
    \phantomsubcaption\label{alg:two-sum}
    \fcolorbox{gray!35}{gray!3}{%
      \begin{minipage}[t]{\dimexpr\linewidth-2\fboxsep-2\fboxrule\relax}
        \centering
        \textbf{(\thesubfigure) \TwoSum{}$(a,b)$}
        \par\smallskip
        \raggedright\small
        \textit{Input:} Floating-point values $a$ and $b$\par
        \textit{Output:} A rounded sum $s$ and correction $t$
        \setlength{\abovedisplayskip}{4pt}
        \setlength{\belowdisplayskip}{0pt}
        \[
          \begin{array}{@{}r@{\hspace{0.6em}}l@{}}
            1: & s \leftarrow \operatorname{RN}(a+b)\\
            2: & a' \leftarrow \operatorname{RN}(s-b)\\
            3: & b' \leftarrow \operatorname{RN}(s-a')\\
            4: & \delta_a \leftarrow \operatorname{RN}(a-a')\\
            5: & \delta_b \leftarrow \operatorname{RN}(b-b')\\
            6: & t \leftarrow \operatorname{RN}(\delta_a+\delta_b)\\
            7: & \mathbf{return}\ (s,t)
          \end{array}
        \]
      \end{minipage}}
  \end{subfigure}

  \vspace{7pt}

  \begin{subfigure}[t]{\columnwidth}
    \phantomsubcaption\label{alg:fast-two-sum}
    \fcolorbox{gray!35}{gray!3}{%
      \begin{minipage}[t]{\dimexpr\linewidth-2\fboxsep-2\fboxrule\relax}
        \centering
        \textbf{(\thesubfigure) \FastTwoSum{}$(a,b)$}
        \par\smallskip
        \raggedright\small
        \textit{Input:} Floating-point values $a$ and $b$\par
        \textit{Output:} A rounded sum $s$ and correction $t$
        \setlength{\abovedisplayskip}{4pt}
        \setlength{\belowdisplayskip}{0pt}
        \[
          \begin{array}{@{}r@{\hspace{0.6em}}l@{}}
            1: & s \leftarrow \operatorname{RN}(a+b)\\
            2: & z \leftarrow \operatorname{RN}(s-a)\\
            3: & t \leftarrow \operatorname{RN}(b-z)\\
            4: & \mathbf{return}\ (s,t)
          \end{array}
        \]
      \end{minipage}}
  \end{subfigure}

  \vspace{7pt}

  \begin{subfigure}[t]{\columnwidth}
    \phantomsubcaption\label{alg:fast-two-mult}
    \fcolorbox{gray!35}{gray!3}{%
      \begin{minipage}[t]{\dimexpr\linewidth-2\fboxsep-2\fboxrule\relax}
        \centering
        \textbf{(\thesubfigure) \FastTwoMult{}$(a,b)$}
        \par\smallskip
        \raggedright\small
        \textit{Input:} Floating-point values $a$ and $b$\par
        \textit{Output:} A rounded product $\pi_h$ and correction $\pi_\ell$
        \setlength{\abovedisplayskip}{4pt}
        \setlength{\belowdisplayskip}{0pt}
        \[
          \begin{array}{@{}r@{\hspace{0.6em}}l@{}}
            1: & \pi_h \leftarrow \operatorname{RN}(ab)\\
            2: & \pi_\ell \leftarrow
                 \operatorname{RN}(ab-\pi_h)
                 \quad\text{(fused multiply-add)}\\
            3: & \mathbf{return}\ (\pi_h,\pi_\ell)
          \end{array}
        \]
      \end{minipage}}
  \end{subfigure}
  \endgroup
\caption{Error-free transformations used by the double-word algorithms
implemented in \FenceDD.
\TwoSum{} and \FastTwoSum{} return a rounded sum and correction;
\FastTwoMult{} returns a rounded product and FMA correction.
\Cref{thm:addition-transformations,thm:fast-two-sum-small,thm:fast-two-mult}
state when each correction is exact.}
  \label{fig:error-free-transformations}
\end{figure}

\subsubsection{Exact corrections for addition.}

Gradual underflow does not affect the corrections
returned by \TwoSum{} and \FastTwoSum{}. When the returned components are
finite, \TwoSum{} returns the rounded sum and its exact rounding error.
\FastTwoSum{} also returns the rounded sum and its exact rounding error
under the usual ordering condition on the operands. 
These results, proved by \citet{Boldo2017}, allow the later
double-word proofs to use the addition transformations without imposing a
lower bound on the magnitude of the operands.

\begin{theorem}[Error-free transformations for addition over IEEE~754]
\label{thm:addition-transformations}
Let $a$ and $b$ be finite values in a radix-two format, and suppose the
returned components are finite.
\begin{enumerate}
  \item If $(s,t)=\TwoSum{}(a,b)$, then
    $s=\operatorname{RN}(a+b)$ and $s+t=a+b$.
  \item If $|b|\leq|a|$ and $(s,t)=\FastTwoSum{}(a,b)$, then $(s,t)$
    satisfies $s=\operatorname{RN}(a+b)$ and $s+t=a+b$.
\end{enumerate}
\end{theorem}

The ordering condition $|b|\leq|a|$ is sufficient but not necessary
for exactness, and it cannot be established for every use of
\FastTwoSum{} in the double-word algorithms. We prove another
sufficient condition needed by the later analysis: assuming finite
returned components, the correction remains exact when
$|a+b|\leq2\nu$.

\begin{theorem}[\FastTwoSum{} for small exact sums]
\label{thm:fast-two-sum-small}
Let $\nu$ be the smallest positive normal value in the format, as defined in
\Cref{eq:flt-thresholds}.  If $(s,t)=\FastTwoSum{}(a,b)$ has finite components
and $|a+b|\leq2\nu$, then $s=\operatorname{RN}(a+b)$ and $s+t=a+b$.
\end{theorem}

The proof appears in Appendix~\ref{app:fast-two-sum-small-proof}; like all 
results in this section, it is checked in Rocq.

\subsubsection{Corrections for multiplication near underflow.}

\FastTwoMult{} returns
$\pi_h=\operatorname{RN}(ab)$ and
$\pi_\ell=\operatorname{RN}(d)$, where $d=ab-\pi_h$ is the exact rounding
error. If $d$ is representable, then $\pi_\ell=d$ and the returned components
sum to $ab$; near underflow, however, $d$ may be too small to represent.
Theorem~2 of \citet[Section~3]{BoldoDaumas2003} proves representability when
$a=n_a2^{e_a}$ and $b=n_b2^{e_b}$ satisfy
$e_a+e_b\geq e_{\min}$. From this result, we derive a sufficient condition on $|ab|$ under which the
returned components sum to the exact product. When this condition does not
hold but both components are finite, we instead bound the difference between
their sum and the exact product.

\begin{theorem}[\FastTwoMult{} over IEEE~754]
\label{thm:fast-two-mult}
Let $a$ and $b$ be finite values in a radix-two format with precision
$p\geq3$ and lower exponent $e_{\min}$. Define
\[
  \eta\triangleq2^{e_{\min}}
  \qquad\text{and}\qquad
  \tau\triangleq2^{e_{\min}+2p-1}.
\]
Suppose $(\pi_h,\pi_\ell)=\FastTwoMult{}(a,b)$ returns two finite components.
Then
\[
  \pi_h=\operatorname{RN}(ab)
  \qquad\text{and}\qquad
  \left|(\pi_h+\pi_\ell)-ab\right|\leq\frac{\eta}{2}.
\]
If $ab=0$ or $|ab|\geq\tau$, then
\[
  \pi_\ell=ab-\pi_h
  \qquad\text{and}\qquad
  \pi_h+\pi_\ell=ab.
\]
\end{theorem}

\begin{proof}[Proof sketch]
Let $d=ab-\pi_h$ be the exact rounding error. Lines~1 and~2 of
\Cref{alg:fast-two-mult} give $\pi_h=\operatorname{RN}(ab)$ and
$\pi_\ell=\operatorname{RN}(d)$, so
$|(\pi_h+\pi_\ell)-ab|=|\pi_\ell-d|$.

If $ab=0$, then $d=\pi_\ell=0$. If $|ab|\geq\tau$, representations
$a=n_a2^{e_a}$ and $b=n_b2^{e_b}$ cannot satisfy
$e_a+e_b<e_{\min}$, since that inequality would imply $|ab|<\tau$.
Theorem~2 of \citet[Section~3]{BoldoDaumas2003} therefore shows that $d$ is
representable, giving $\pi_\ell=d$ and exact reconstruction.

In the remaining case, $0<|ab|<\tau$. The FLT rounding bound gives
$|d|<2\nu$, where $\nu=2^{e_{\min}+p-1}$. Representable values in this range
are spaced by $\eta$, so $|\pi_\ell-d|\leq\eta/2$, which proves the stated
error bound.
\end{proof}

\subsection{Numerical Guarantee for Division of a Double Word by a
Floating-Point Value}
\label{sec:division-numerics}

Division of a double word by a floating-point value shows how the preceding
error-free transformation results combine within a complete double-word
operation. We identify sufficient conditions under which the returned pair
forms a double word and its represented value satisfies an error bound with
relative and absolute terms. 
The proof accounts for both effects of the finite exponent
range: the relevant computations must remain finite, and operations involving
subnormal values may require direct analysis in FLT.

Given a double word $x=x_h+x_\ell$ and a nonzero floating-point divisor $y$,
the exact quotient is
\begin{equation}
  q=\frac{x_h+x_\ell}{y}.
  \label{eq:dw-div-exact-quotient}
\end{equation}
The functional model \DWDivFP{} represents the rounded operations used to
approximate this quotient. \Cref{alg:dw-div-fp} gives its mathematical
definition. Each $\operatorname{RN}$ denotes rounding to nearest with ties to even.

\begin{figure}[t]
  \centering
  \begingroup
  \setlength{\fboxsep}{6pt}
  \setlength{\fboxrule}{0.4pt}
  \fcolorbox{gray!35}{gray!3}{%
    \begin{minipage}[t]{\dimexpr\linewidth-2\fboxsep-2\fboxrule\relax}
      \centering
      \textbf{\DWDivFP{}$(x_h,x_\ell,y)$}
      \par\smallskip
      \raggedright\small
      \textit{Input:} A double word $(x_h,x_\ell)$ and a nonzero
      floating-point value $y$\par
      \textit{Output:} A pair $(z_h,z_\ell)$ approximating $q$
      \setlength{\abovedisplayskip}{4pt}
      \setlength{\belowdisplayskip}{0pt}
      \[
        \begin{array}{@{}r@{\hspace{0.6em}}l@{}}
          1: & t_h \leftarrow \operatorname{RN}(x_h/y)\\
          2: & (\pi_h,\pi_\ell) \leftarrow \FastTwoMult{}(t_h,y)\\
          3: & \delta_h \leftarrow \operatorname{RN}(x_h-\pi_h)\\
          4: & \delta_t \leftarrow \operatorname{RN}(\delta_h-\pi_\ell)\\
          5: & \delta \leftarrow \operatorname{RN}(\delta_t+x_\ell)\\
          6: & t_\ell \leftarrow \operatorname{RN}(\delta/y)\\
          7: & (z_h,z_\ell) \leftarrow \FastTwoSum{}(t_h,t_\ell)\\
          8: & \mathbf{return}\ (z_h,z_\ell)
        \end{array}
      \]
    \end{minipage}}
  \endgroup
  \caption{The functional model \DWDivFP{} for dividing a double word by a
  floating-point value.}
  \label{alg:dw-div-fp}
\end{figure}

The model follows the identity
\[
  q=\frac{x}{y}
    =t_h+\frac{x-t_hy}{y}.
\]
Line~1 computes the leading quotient $t_h$.
Lines~2--5 approximate the residual $x-t_hy$, including the contribution of
the low word $x_\ell$. Line~6 divides this residual by $y$ to obtain the
correction $t_\ell$. Line~7 constructs the final pair.

\subsubsection{Sufficient Conditions for the Numerical Guarantee}
\label{sec:division-conditions}

We state sufficient conditions for \DWDivFP{} to return a finite
double word satisfying an error bound. We state the operand
conditions separately from the assumption that the returned
components are finite, then give stronger conditions under which
result finiteness follows.

For operands $x_h,x_\ell,y$, let
$(z_h,z_\ell)=\DWDivFP{}(x_h,x_\ell,y)$, and let $t_h$ be the rounded leading
quotient computed on Line~1 of \Cref{alg:dw-div-fp}. When $y\neq0$, write
$q_0=x_h/y$ for the exact leading quotient. We write
$\operatorname{Fin}(v_1,\ldots,v_n)$ when all its arguments are finite. Let
$\Omega$ be the least power of two greater than every finite magnitude in the
format; for binary64, $\Omega=2^{1024}$.

\begin{definition}[Conditions for \DWDivFP{}]
\label[definition]{def:dw-div-conditions}
Define
\begin{equation}
\begin{aligned}
\mathcal{I}_{\mathrm{div}}(x_h,x_\ell,y)
  \quad\triangleq\quad&
  (x_h,x_\ell)\text{ is a double word}\\
  &{}\land
  \left(x_h=0
    \;\lor\;
    2^{e_{\min}+2p}\leq |x_h|\right),\\[2pt]
\mathcal{F}_{\mathrm{div}}(x_h,x_\ell,y)
  \quad\triangleq\quad&
  \operatorname{Fin}(z_h,z_\ell),\\[2pt]
\mathcal{S}_{\mathrm{div}}(x_h,x_\ell,y)
  \quad\triangleq\quad&
  \mathcal{I}_{\mathrm{div}}(x_h,x_\ell,y)
  \land \operatorname{Fin}(x_h,x_\ell,y,t_h)\\
  &{}\land |x_h|\leq\Omega/4
  \land |q_0|\leq\Omega/4.
\end{aligned}
\label{eq:dw-div-conditions}
\end{equation}
\end{definition}

The condition $\mathcal{I}_{\mathrm{div}}$ requires the dividend to be a
double word and imposes a lower bound on its high word unless the dividend is
zero. It places no lower bound on $q_0$, so the leading quotient may be
subnormal. The condition $\mathcal{F}_{\mathrm{div}}$ states that both
returned words are finite, as required for the error bound and double-word
relation. The stronger condition $\mathcal{S}_{\mathrm{div}}$ adds finiteness
and upper bounds sufficient to derive
$\mathcal{F}_{\mathrm{div}}$ without assuming it separately.

\begin{lemma}[Finiteness from the Stronger Condition]
\label[lemma]{lem:dw-div-safe-implies-finite}
If $\mathcal{S}_{\mathrm{div}}(x_h,x_\ell,y)$ holds, then
$\mathcal{F}_{\mathrm{div}}(x_h,x_\ell,y)$ holds.
\end{lemma}

\subsubsection{Error Bounds and Double-Word Result}

The main theorem assumes
$\mathcal{I}_{\mathrm{div}}$ and $\mathcal{F}_{\mathrm{div}}$. Under these
conditions, the functional model returns a double word and satisfies an error
bound with relative and absolute terms. Under additional magnitude conditions,
the absolute term is unnecessary. By the preceding finiteness lemma, the same
conclusions also follow from $\mathcal{S}_{\mathrm{div}}$.

\begin{theorem}[Numerical guarantee for \DWDivFP{}]
\label{thm:dw-div-fp-numerical}
Let the underlying format be radix two with precision $p\geq4$, lower
exponent $e_{\min}$, unit roundoff $u=2^{-p}$, and
$\eta=2^{e_{\min}}$. Let $\nu=2^{e_{\min}+p-1}$ be its smallest positive
normal value, and let
$(z_h,z_\ell)=\DWDivFP{}(x_h,x_\ell,y)$.
Suppose
$\mathcal{I}_{\mathrm{div}}(x_h,x_\ell,y)$ and
$\mathcal{F}_{\mathrm{div}}(x_h,x_\ell,y)$ hold. These conditions imply
$y\neq0$. Let $q$ be the exact quotient defined in
\Cref{eq:dw-div-exact-quotient}, and let $q_0=x_h/y$ be the exact leading
quotient. Then $(z_h,z_\ell)$ is a double word:
\[
  z_h=\operatorname{RN}(z_h+z_\ell).
\]
Its represented value satisfies
\begin{equation}
  \left|(z_h+z_\ell)-q\right|
    \leq 3u^2|q|+\eta.
  \label{eq:dw-div-fp-mixed}
\end{equation}
Let $\delta$ denote the residual computed on Line~5 of
\Cref{alg:dw-div-fp}. If $|q_0|\geq\nu$ and either $\delta=0$ or
$|\delta/y|\geq\nu$, the sharper bound holds:
\begin{equation}
  \left|(z_h+z_\ell)-q\right|
    \leq 3u^2|q|.
  \label{eq:dw-div-fp-relative}
\end{equation}
\end{theorem}

\begin{proof}[Proof sketch]
The conditions imply that $y\neq0$ and that the relevant intermediate results
are finite. If $x_h=0$, the double-word condition gives $x_\ell=0$, and the
result follows directly. Otherwise,
$|x_h|\geq2^{e_{\min}+2p}$.

First suppose $|q_0|\geq\nu$. Rounding on Line~1 agrees in FLT
and FLX. The lower bound on
$|x_h|$ gives $|t_hy|\geq2^{e_{\min}+2p-1}$, so
\Cref{thm:fast-two-mult} gives $\pi_h+\pi_\ell=t_hy$. Sterbenz's lemma and
representability then carry the FLX residual analysis through Lines~3--5
\citep{Sterbenz1974}. If $\delta=0$ or $|\delta/y|\geq\nu$, Line~6 also agrees
in the two formats and gives \Cref{eq:dw-div-fp-relative}; otherwise, its FLT
and FLX results differ by at most $\eta$, giving the mixed bound. Finally,
$|t_\ell|\leq|t_h|$, so \FastTwoSum{} preserves their sum and returns a double
word by \Cref{thm:addition-transformations}.

Now suppose $|q_0|<\nu$. The fixed spacing in FLT gives
$|t_h-q_0|\leq\eta/2$. The double-word
condition gives $|x_\ell/y|<\eta/2$, and hence $|t_h-q|\leq\eta$. The lower
bound on $|x_h|$ also implies $|y|>2^{p+1}$. Applying the FLT rounding bounds
and \Cref{thm:fast-two-mult} to Lines~2--5 gives
\begin{equation}
  |\delta-(x_h+x_\ell-t_hy)|
    \leq\frac{\eta|y|}{2},
  \qquad
  \left|\frac{\delta}{y}-(q-t_h)\right|\leq\frac{\eta}{2}.
  \label{eq:dw-div-subnormal-residual}
\end{equation}
Line~6 adds at most another $\eta/2$, and $|t_h+t_\ell|\leq2\nu$.
\Cref{thm:fast-two-sum-small} therefore shows that Line~7 preserves the sum
and returns a double word, with
\[
  |(z_h+z_\ell)-q|
    \leq\eta
    \leq3u^2|q|+\eta.
\]
This proves both error bounds and the double-word property.
\end{proof}

By \Cref{lem:dw-div-safe-implies-finite},
$\mathcal{S}_{\mathrm{div}}$ implies $\mathcal{F}_{\mathrm{div}}$ and already
includes $\mathcal{I}_{\mathrm{div}}$. Consequently,
\Cref{thm:dw-div-fp-numerical} gives the following two alternatives.

\begin{corollary}[Conditions for the Division Guarantee]
\label[corollary]{cor:dw-div-conditions}
The conclusions of \Cref{thm:dw-div-fp-numerical} hold if
\[
  \mathcal{I}_{\mathrm{div}}(x_h,x_\ell,y)
  \land
  \mathcal{F}_{\mathrm{div}}(x_h,x_\ell,y)
\]
holds, or if $\mathcal{S}_{\mathrm{div}}(x_h,x_\ell,y)$ holds.
\end{corollary}

\Cref{thm:division-checks} proves that the weak and result finiteness checks in
the default wrapper establish the first alternative, while the strong check in
the alternative wrapper establishes the second.

\subsection{Numerical Guarantees for the Remaining Operations}
\label{sec:remaining-numerical-conditions}

The five remaining functional models also return double words and satisfy the
error bounds summarized in \Cref{tab:operation-preview} under
operation-specific conditions. Rather than repeat their algorithms and full
proofs, we collect the conditions needed to extend the existing analyses to
the finite IEEE~754 exponent range and state one common theorem.

Write $x=(x_h,x_\ell)$ and, where present, $y=(y_h,y_\ell)$ for the
double-word operands, and let $a$ denote a floating-point operand. For
addition and multiplication, the returned components must be finite, but no
additional condition on the operands is required beyond the double-word
assumptions. Division of two double words additionally requires
%
\begin{equation}
  2^{e_{\min}+2p}
  \leq |y_h|
  \leq \frac{\Omega}{2^{p+5}}.
  \label{eq:dw-div-dw-divisor-condition}
\end{equation}
%
In every case, finiteness may either be assumed directly for the returned
components or established from the sufficient conditions in
\Cref{tab:remaining-numerical-conditions}.
%
\begin{table}[t]
\centering
\caption{Sufficient conditions for the five remaining functional
models to return finite components. Each row provides an
alternative to the result-finiteness assumption in
\Cref{thm:remaining-numerical-guarantees}; all other hypotheses
still apply. Here $x=(x_h,x_\ell)$ and $y=(y_h,y_\ell)$ are
double-word operands, $a$ is a floating-point operand, and
$\operatorname{Fin}$ states that all listed values are finite.
The constant $\Omega$ is the least power of two greater than
every finite magnitude in the underlying format.}
\label{tab:remaining-numerical-conditions}
\setlength{\tabcolsep}{3pt}
\begin{tabular}{@{}p{0.31\columnwidth}p{0.64\columnwidth}@{}}
\toprule
Operation & Sufficient conditions for a finite result \\
\midrule
Addition $(x,a)$
  & $\operatorname{Fin}(x_h,x_\ell,a)$;
    $|x_h|,|a|\leq\Omega/4$ \\
\addlinespace
Addition $(x,y)$
  & $\operatorname{Fin}(x_h,x_\ell,y_h,y_\ell)$;
    $|x_h|,|y_h|\leq\Omega/4$ \\
\addlinespace
Multiplication $(x,a)$
  & $\operatorname{Fin}(x_\ell,\operatorname{RN}(x_h a))$;
    $|x_h a|\leq\Omega/2$ \\
\addlinespace
Multiplication $(x,y)$
  & $\operatorname{Fin}(x_\ell,y_\ell,\operatorname{RN}(x_h y_h))$;
    $|x_h y_h|\leq\Omega/2$ \\
\addlinespace
Division $(x,y)$
  & $\operatorname{Fin}(x_h,x_\ell,y_h,y_\ell)$;
    $|x_h/y_h|\leq\Omega/4$ \\
\bottomrule
\end{tabular}
\end{table}
%
Under these conditions, each functional model returns a double word and
satisfies its operation-specific error bound:
%
\begin{theorem}[Numerical guarantees for the remaining operations]
\label{thm:remaining-numerical-guarantees}
Let the underlying format be radix two with precision $p$, unit
roundoff $u=2^{-p}$, and smallest positive representable value
$\eta=2^{e_{\min}}$. Choose one of the five operations in
\Cref{tab:remaining-numerical-conditions}, and let $(z_h,z_\ell)$
be the pair returned by the corresponding functional model.
Let $r$ be the exact real result of the operation on the values
represented by the operands.

Assume $p\geq3$ for the addition operations and multiplication by a
floating-point value, $p\geq5$ for multiplication of two double
words, and $p\geq14$ for division of two double words. Suppose every
double-word operand satisfies \Cref{def:double-word}. For division
of two double words, also suppose
\Cref{eq:dw-div-dw-divisor-condition} holds.

Finally, suppose either that the returned components are finite,
$\operatorname{Fin}(z_h,z_\ell)$, or that the sufficient conditions
for finiteness in the corresponding row of
\Cref{tab:remaining-numerical-conditions} hold.

Then both returned components are finite and form a double word:
\[
  z_h=\operatorname{RN}(z_h+z_\ell).
\]
The represented value $\widehat r=z_h+z_\ell$ satisfies the
corresponding error bound on $|\widehat r-r|$ in
\Cref{tab:operation-preview}, using $u$ and $\eta$ for the underlying
format.
\end{theorem}

\paragraph{Proof overview.}
\emph{Addition.}
The error-free transformations for addition continue to return the exact
rounding error when intermediate results are subnormal by
\Cref{thm:addition-transformations}. Consequently, the existing relative
error bounds require no lower bound on the operand magnitudes.

\emph{Multiplication.}
When a product correction is too small to be represented exactly,
\Cref{thm:fast-two-mult} bounds the lost correction. Propagating that error
through the remaining computation produces the absolute terms in the 
\Cref{tab:operation-preview} bounds.

\emph{Division of two double words.}
In contrast to division by a floating-point value in
\Cref{alg:dw-div-fp}, division of two double words begins by approximating the
reciprocal of $y_h$. The lower bound in
\Cref{eq:dw-div-dw-divisor-condition} prevents this approximation from
overflowing. The upper bound implies $\eta|y_h|\leq u^2/4$, allowing the error
from a lost product correction to be absorbed into the relative term. 
The bound on the quotient in \Cref{tab:remaining-numerical-conditions} limits the
magnitude of the leading quotient and is sufficient to prove that the returned
components are finite.

\section{Verification of \FenceDD}
\label{sec:mechanization}

The numerical theorems in \Cref{sec:numerics} concern the functional models
defined there, whereas users call the public C wrappers in
\Cref{sec:overview}. Following
earlier layered verifications of numerical C programs
\citep{AppelBertot2020,Kellison2023LAProof,Tekriwal2023Jacobi}, we prove three
obligations: kernel refinement, soundness of the runtime checks, and their
composition into each wrapper's public contract.  We illustrate these proofs
using division of a double word by a binary64 value; the other operations
follow the same structure.

\subsection{C Kernels Implement Their Functional Models}
\label{sec:kernel-refinement}

The numerical theorem in \Cref{thm:dw-div-fp-numerical} concerns the
functional model \DWDivFP{}, while the wrappers call the C kernel in
\Cref{fig:dw-div-fp-kernel}.  Suppose that \texttt{x} initially stores
$(x_h,x_\ell)$ and that \texttt{dw\_div\_fp\_kernel} returns with
\texttt{z} storing $(z_h,z_\ell)$.  Using VST, we prove
\begin{equation}
  (z_h,z_\ell)=\DWDivFP(x_h,x_\ell,y),
  \label{eq:dw-div-fp-kernel-refinement}
\end{equation}
where \DWDivFP{} is instantiated with binary64. The Rocq definition of
\DWDivFP{} uses the IEEE~754 types and operations provided by VCFloat2
\citep{AppelKellison2024}. In the VST proof, \texttt{FPCompCert} bridge
lemmas relate these operations to CompCert, while the specifications of the
transformation kernels account for the internal C calls.  This refinement identifies
the computed pair but requires no numerical condition and proves no error
bound.  We prove an analogous result for every arithmetic kernel in the
library.

\subsection{Runtime Checks Establish the Mathematical Conditions}
\label{sec:check-soundness}

To apply the numerical theorem to either wrapper for dividing a double-double
by a binary64 value, we must show that successful runtime checks establish its
mathematical hypotheses. The
following theorem connects the weak and strong checks in
\Cref{fig:dw-div-fp-input-check,fig:dw-div-fp-check} and the finiteness check
in \Cref{fig:finite-checks} to the conditions in
\Cref{sec:division-conditions}.

\begin{theorem}[Soundness of the Checks for Division by a Binary64 Value]
\label{thm:division-checks}
Suppose that \texttt{x} stores the binary64 pair $(x_h,x_\ell)$.

\begin{enumerate}
  \item If the weak check returns $1$, then
    $\mathcal{I}_{\mathrm{div}}(x_h,x_\ell,y)$ holds.

  \item If the strong check returns $1$, then
    $\mathcal{S}_{\mathrm{div}}(x_h,x_\ell,y)$ holds.
\end{enumerate}

If \texttt{z} stores $(z_h,z_\ell)$ and the finiteness check returns $1$,
then both components are finite.
\end{theorem}

The proofs account for short-circuit control flow using the specifications of
the finiteness functions, \texttt{fabs}, and \texttt{ldexp}.  The main
arithmetic argument derives the required bound on the exact quotient from the
rounded quotient comparison, using exact threshold representability and
monotonicity of rounding.  These results prove soundness, not completeness: a
failed runtime check need not imply that the corresponding mathematical
condition is false.

\subsection{Verified Contracts for Dividing a Double-Double by a Binary64
Value}
\label{sec:division-wrapper-contracts}

The wrapper proofs combine kernel refinement, check soundness, and the
numerical theorem.  \Cref{fig:dw-div-fp-vst-spec} shows the clauses of the VST
specification for the default wrapper in
\Cref{fig:default-division-wrapper} that state its numerical contract; routine
memory assertions are omitted.

\begin{figure}[t]
  \begin{lstlisting}[
    language=Coq,
    basicstyle=\ttfamily\normalsize,
    numbers=left
  ]
Definition dw_div_fp_spec :=
  DECLARE _dw_div_fp
  WITH z_ptr: val, x_ptr: val,
       x_pair: ftype Tdouble *
               ftype Tdouble,
       y: ftype Tdouble
  PRE [tptr dword, tptr dword, tdouble]
    PROP()
    ...
  POST [tint]
    EX status : int,
    PROP(default_wrapper_status
           status;
      status = Int.zero ->
      DWDivFP_guarantee
        (fst x_pair)
        (snd x_pair) y)
    RETURN(Vint status)
    SEP((if Int.eq status Int.zero
      then data_at Tsh dword
        (f2val
          (DWDivFP
            (fst x_pair)
            (snd x_pair) y))
        z_ptr
      else data_at_ Tsh dword z_ptr);
      ...).
  \end{lstlisting}
\caption{Relevant clauses of the VST specification for the default division
wrapper. On success, the postcondition supplies the numerical guarantee.
Omitted assertions describe the parameters and unchanged input object.}
  \label{fig:dw-div-fp-vst-spec}
\end{figure}


The empty \texttt{PROP} clause means that the caller supplies no numerical
hypothesis.  On normal return, the postcondition states that the output object
stores the binary64 result of \DWDivFP{} and that
\texttt{DWDivFP\_guarantee} holds.  In the Rocq
development, \texttt{DWDivFP\_guarantee} names the conclusion of
\Cref{thm:dw-div-fp-numerical}: the represented result satisfies the stated
error bound, and the returned pair satisfies the double-word relation.

For the default wrapper, the weak check establishes
$\mathcal{I}_{\mathrm{div}}$, kernel refinement identifies the result with
\DWDivFP{}, and the result check establishes $\mathcal{F}_{\mathrm{div}}$.
For the alternative wrapper, the strong check establishes
$\mathcal{S}_{\mathrm{div}}$ and hence result finiteness.  These are the two
alternatives in \Cref{cor:dw-div-conditions}, so both wrappers provide the same
contract.

\subsection{Compilation and the Verification Boundary}
\label{sec:verification-boundary}

The C proofs assume that \texttt{double} is binary64 and that arithmetic uses
round-to-nearest, ties-to-even, with gradual underflow.
Compilation and execution must preserve these semantics without reassociating
floating-point expressions or flushing subnormal values to zero.
Floating-point exceptions are nontrapping, which allows the default wrapper to
reach the result check on Line~9 of \Cref{fig:default-division-wrapper} when
the kernel produces an infinity or NaN \citep{IEEE754-2019}.
The verified functions call \texttt{fma}, \texttt{fabs}, \texttt{ldexp}, and
\texttt{exit}.  In the proofs, we use the behavior described by their VST
specifications, but the platform implementations of these functions and the
floating-point hardware are outside the verification boundary
\citep{Appel2023VSTlib}.

We prove the C-level contracts using VST.  For a CompCert build, we apply
semantic preservation to carry these results to assembly
\citep{Leroy2009,Boldo2013CompCert}.  Compilation with another
compiler is not covered by this result and must independently preserve the
required IEEE~754 semantics.

\section{Evaluation}
\label{sec:evaluation}

The verified contracts establish numerical guarantees on normal return.  We
evaluate three complementary empirical questions: the
tightness of the proved error bounds, the practical significance of calls
rejected by the runtime checks, and the performance cost of the
wrappers.  We investigate the following research questions:

\begin{description}
  \item[(RQ1) Usefulness of the error bounds.]
    For results produced by accepted calls to the default wrappers for the six
    double-word operations, how closely do the observed errors approach the
    proved bounds?
    As secondary context, how do the errors produced by QD compare on the same
    inputs?

  \item[(RQ2) Usefulness of the wrappers.]
    For generated inputs that cause a runtime check in each wrapper to
    fail, how often does QD return a finite, valid double word whose
    MPFR-measured error is no greater than the corresponding \FenceDD{} error
    bound?

  \item[(RQ3) Cost of the wrappers.]
    On a common set of inputs for which every evaluated call is accepted, what
    runtime overhead does each wrapper add relative to the
    corresponding unchecked kernel, and how does its end-to-end runtime
    compare with QD?
\end{description}

\subsection{Experimental Setup}
\label{sec:evaluation-setup}

We compare \FenceDD{} with QD \citep{Hida2001,QD}, using MPFR to compute
reference values at greater precision \citep{Fousse2007}.  We evaluate both
wrappers for each of the six double-word operations and the single public
wrapper for each error-free transformation.  For every test, both libraries
receive the same binary64 operands, including the same high and low components
of each double-word operand.

To exercise the binary64 range, we generate the same number of inputs at each
base-two exponent.  Binary64 operands cover exponents from $-1074$ through
$1023$, while double-word inputs cover exponents from $-969$ through $1022$.
For 95\% of the double-word inputs, MPFR constructs a 106-bit source value,
which is rounded into high and low components.  The remaining 5\% have a
binary64 high component and a zero low component.  The narrower exponent range
keeps all 106 source bits representable while ensuring that the high component
is finite.  Signs are balanced, and every generated double-word operand is
valid according to \Cref{def:double-word}.

We classify a call as accepted when all runtime checks succeed and the
function returns a result.  A call is rejected when a runtime check stops its
execution.  RQ1 and RQ2 use \RQOneSamplesPerExponent{} inputs at each exponent,
giving 79,680 inputs per double-word operation and 83,920 per transformation.
All floating-point computations use rounding to nearest with ties to even and
gradual underflow.  \Cref{app:evaluation-details} gives the complete input
construction, reference computation, boundary tests, and timing protocol.

\subsection{Usefulness of the Error Bounds (RQ1)}

To address RQ1, we evaluate the six double-word operations on generated inputs
for which calls to the default wrapper are accepted.  MPFR reconstructs the real
values represented by the input and output pairs and computes the
corresponding reference operation at 2112-bit precision.  If $r$ is the MPFR
reference and $\widehat r$ is the value represented by the returned pair, the
observed error is $|\widehat r-r|$.

The bounds in \Cref{tab:operation-preview} have the form
\begin{equation}
  \begin{aligned}
    B(r) &= C u^2|r|+K\eta,\\
    B_{\mathrm{rel}}(r) &= C u^2|r|,
    & B_{\mathrm{abs}} &= K\eta.
  \end{aligned}
  \label{eq:rq1-error-bound}
\end{equation}
We refer to $B_{\mathrm{rel}}(r)$ and $B_{\mathrm{abs}}$ as the relative and
absolute terms.  An input lies in the relative region
when $B_{\mathrm{rel}}(r)\geq B_{\mathrm{abs}}$ and in the absolute region
otherwise.  The addition bounds have no absolute term, so all nonzero addition
results lie in the relative region.

For every result with $B(r)>0$, we measure how closely its error approaches
the bound using
\begin{equation}
  \rho=\frac{|\widehat r-r|}{B(r)}.
  \label{eq:rq1-error-ratio}
\end{equation}
A value of $\rho$ near one indicates that the observed error approaches the
proved bound.  Smaller values indicate that the bound is less tightly attained
on that input.  The numerical theorem implies $\rho\leq1$ for every
\FenceDD{} result produced by an accepted call.  Addition results for which $B(r)=0$ are
reported separately and must be exact.

The QD routines may use a different sequence of binary64 operations from the
corresponding \FenceDD{} kernels.  To provide a common numerical scale, we
divide each finite QD error by the corresponding \FenceDD{} error bound
$B(r)$.  This bound is not proved for QD, and the QD and \FenceDD{} ratios are
not expected to be equal.  Nonfinite QD results are reported separately.

\begin{table}[t]
  \caption{RQ1: Mean and maximum values of the error ratio $\rho$ from
  \Cref{eq:rq1-error-ratio}, separated into the regions determined by the terms
  of the \FenceDD{} error bound in \Cref{eq:rq1-error-bound}.  Region share is
  the percentage of results produced by accepted \FenceDD{} calls that fall in
  each region; the shares for an operation sum to 100\%.  For QD,
  the value of the \FenceDD{} error bound provides a common scale; the bound is
  not proved for QD.  Region shares are reported to two decimal places and
  ratios to four decimal places.}
  \label{tab:rq1-tightness}
  \centering
  \setlength{\tabcolsep}{2pt}
  \begin{tabular}{@{}llllll@{}}
    \toprule
    & & \multicolumn{2}{c}{\FenceDD{} $\rho$}
        & \multicolumn{2}{c}{QD $\rho$} \\
    \cmidrule(lr){3-4}\cmidrule(l){5-6}
    Operation & Share
              & Mean
              & Max
              & Mean
              & Max \\
    \midrule
    \multicolumn{6}{@{}l}{\emph{Relative region}} \\[2pt]
    \input{generated/rq1_relative_compact_rows}\\
    \addlinespace
    \multicolumn{6}{@{}l}{\emph{Absolute region}} \\[2pt]
    \input{generated/rq1_absolute_compact_rows}\\
    \bottomrule
  \end{tabular}
\end{table}

\paragraph{Answer to RQ1.}
\Cref{tab:rq1-tightness} shows that the proved bounds can be approached, but
are conservative for most generated inputs.  Across
\RQOneNormalReturnCount{} results produced by accepted calls, no \FenceDD{}
error exceeded its proved bound.  The largest observed values of $\rho$ were
\RQOneRelativeMaximumRatio{} for \RQOneRelativeMaximumOperation{} in the
relative region and \RQOneAbsoluteMaximumRatio{} for
\RQOneAbsoluteMaximumOperation{} in the absolute region.  By contrast, the
mean for every operation was at most \RQOneRelativeMaximumMean{} in the
relative region.  For every operation with an absolute term, the mean in the
absolute region was at most \RQOneAbsoluteMaximumMean{}.  The absolute region
contained 14--16\% of the results for these operations, so these measurements
represent a substantial portion of the generated tests rather than isolated
boundary cases.
Thus, some observed errors came close to the proved limits, while the bounds
retained substantial margin on average.

As secondary context, QD produced the same observed ratios for both addition
operations and \DWDivFP{}.  Its relative errors were somewhat larger for the
two multiplication operations and smaller for \textsc{DWDivDW}.  One
\DWTimesFP{} result reached $\rho=1.038$, exceeding $B(r)$ by
\RQOneQDExcessPercent{}.  QD may
compute the same mathematical operation using a different algorithm, so
$B(r)$ is not a proved bound for its result.  We use $B(r)$ only to place the
observed QD and \FenceDD{} errors on a common scale.

\FloatBarrier
\subsection{Usefulness of the Wrappers (RQ2)}

A call is rejected when a runtime check cannot establish a condition required
by its verified contract.  To assess the practical significance of these
rejections, we run the corresponding QD operation on every generated input
associated with a rejected \FenceDD{} call and compare its result with an MPFR
reference.

For the six double-word operations, we consider separately calls rejected by
runtime checks in the default wrappers and calls rejected by runtime checks in
the strong wrappers.  These populations overlap: every call rejected in a
default wrapper is also rejected in the corresponding strong wrapper.
Percentages are therefore calculated separately using the number of rejected
calls in the indicated population.  Because the generated tests contain the
same number of inputs at each exponent, these percentages describe our tests
and do not estimate rejection rates in application workloads.

For a double-word operation, the numerical criterion requires QD to return a
finite, valid double word whose error is no greater than the corresponding
\FenceDD{} bound $B(r)$ from \Cref{eq:rq1-error-bound}.  The bound is not proved
for QD; we use it only as a common numerical criterion.  No finite QD result in
these populations was invalid or exceeded $B(r)$, so the table omits those
zero categories.

\begin{table}[t]
  \caption{RQ2: QD results for inputs on which a runtime check stopped a call to
  the indicated wrapper.  Rejected calls gives the population size; the other
  entries are percentages of that population.  The default and strong
  populations overlap.  No finite QD result was invalid or exceeded $B(r)$, so
  those zero categories are omitted.}
  \label{tab:rq2-double-word}
  \centering
  \setlength{\tabcolsep}{3pt}
  \begin{tabular}{@{}llrrr@{}}
    \toprule
    Operation & Wrapper & \shortstack{Rejected\\calls}
              & \shortstack{QD\\nonfinite}
              & \shortstack{Valid,\\error $\leq B(r)$} \\
    \midrule
    \input{generated/rq2_double_word_compact_rows}\\
    \bottomrule
\end{tabular}
\end{table}

Calls stopped by runtime checks usually occurred on inputs for which QD also
failed the numerical criterion.  QD returned a finite, valid double word with
error no greater than the corresponding \FenceDD{} bound for only
\RQTwoDefaultQDWithinPercent{} of calls stopped in the default wrappers and
\RQTwoStrongQDWithinPercent{} of calls stopped in the strong wrappers.

The operation-level results expose important differences.  Every QD result
associated with a rejected call to a default multiplication wrapper was
nonfinite.  Among calls rejected by runtime checks in the default wrappers, QD
met the numerical criterion only for division.  By contrast, QD met the
criterion for every call rejected by the strong addition checks, showing that
these sufficient checks can be conservative.

Each error-free transformation has one public wrapper; it is not classified as
default or strong.  For inputs associated with rejected calls, we compare
\TwoSum{}, \FastTwoSum{}, and \FastTwoMult{} with QD's
\texttt{two\_sum}, \texttt{quick\_two\_sum}, and \texttt{two\_prod},
respectively.  Because the verified contracts require exact reconstruction, a
finite, valid QD pair either reconstructs the MPFR result exactly or fails the
criterion.

\begin{table}[t]
  \caption{RQ2: QD results for inputs on which a runtime check stopped a call to
  the public wrapper for each error-free transformation.  Rejected calls gives
  the population size; the other entries are percentages of that population.}
  \label{tab:rq2-transformations}
  \centering
  \setlength{\tabcolsep}{2pt}
  \begin{tabular}{@{}lrrrrr@{}}
    \toprule
    Transformation & \shortstack{Rejected\\calls}
                   & Nonfinite
                   & Invalid
                   & \shortstack{Not\\exact}
                   & Exact \\
    \midrule
    \input{generated/rq2_eft_compact_rows}\\
    \bottomrule
  \end{tabular}
\end{table}

The results reinforce the pattern in \Cref{tab:rq2-double-word}.  No generated
call to the \TwoSum{} wrapper was rejected.  On inputs associated with rejected
calls to the wrappers for
\FastTwoSum{} and \FastTwoMult{}, QD reconstructed the exact result in only
\RQTwoFastSumExactPercent{} and \RQTwoFastMultExactPercent{} of cases,
respectively.

\paragraph{Where the two double-word wrappers differ.}
For \RQTwoDefaultOnly{} inputs, every runtime check in the default wrapper
succeeded, while the strong check in the corresponding strong wrapper failed.
For multiplication, these cases had a leading-product exponent of $1023$; for
division, they had leading-quotient exponents of $1022$ or $1023$.  The
addition wrappers differed on 196 inputs because the strong checks bound the
operands separately even though the computed sum remained finite.  Checking
result finiteness after kernel execution therefore allows calls to the default
wrapper to be accepted when the corresponding calls to the strong wrapper are
rejected.

\paragraph{Answer to RQ2.}
On these generated tests, runtime checks usually stopped calls on inputs for
which QD also failed the corresponding numerical criterion.  The
operation-level results nevertheless confirm that the checks are sufficient
rather than necessary, particularly for addition.  Because QD may use
different algorithms, these measurements do not establish how an unchecked
\FenceDD{} kernel would behave on the same inputs.

\FloatBarrier
\subsection{Performance Cost of the Wrappers (RQ3)}

To address RQ3, we time the default and strong wrappers for the six
double-word operations and the wrapper for each of the three transformations.
Each timing test uses \RQThreeInputCount{} inputs
whose operand exponents lie between $-100$ and $100$, together with
\RQThreeTrialCount{} trials.  All runtime checks succeed for these inputs.
Before timing, we confirm that every wrapper produces the same
result as its unchecked kernel.

For the comparison with the kernels, each wrapper and its kernel are compiled
with CompCert.  For the comparison with QD, both libraries are compiled with
Clang.  The latter build is used only for the performance comparison and is
not covered by the CompCert correctness theorem.  The experiments run on an
Apple M4 Max with 16 cores.  Compiler versions, options, and the complete
timing procedure appear in \Cref{app:evaluation-timing}.

\Cref{tab:rq3-performance} summarizes the resulting timing measurements for
the default double-word wrappers and the transformation wrappers.  The
underlying times and results for the strong wrappers appear in
\cref{tab:rq3-performance-detail}.

\begin{table}[t]
  \caption{RQ3: Performance of the default double-word wrappers and the
  transformation wrappers.  Under CompCert, Wrapper/Kernel is the ratio to the
  unchecked kernel and Added is the paired difference in nanoseconds per call.
  Under Clang, Wrapper/QD is the ratio to QD.  Entries are medians
  [interquartile ranges] over paired trials.}
  \label{tab:rq3-performance}
  \centering
  \setlength{\tabcolsep}{3pt}
  \begin{tabular}{@{}lrrr@{}}
    \toprule
    & \multicolumn{2}{c}{CompCert} & Clang \\
    \cmidrule(lr){2-3}\cmidrule(l){4-4}
    Operation & Wrapper/Kernel & Added (ns) & Wrapper/QD \\
    \midrule
    \input{generated/rq3_compact_rows}\\
    \bottomrule
  \end{tabular}
\end{table}

We rotate the measurement order within paired trials to limit ordering effects
and compute differences and ratios from the paired measurements.

\paragraph{Answer to RQ3.}
\Cref{tab:rq3-performance} and \Cref{tab:rq3-performance-detail} show that the
wrappers have a noticeable relative cost because their kernels are small, but
add only a few nanoseconds per call.  Under CompCert, the default wrappers took
\RQThreeDefaultMinimumKernelRatio{}--\RQThreeDefaultMaximumKernelRatio{} times
as long as their kernels and added
\RQThreeDefaultMinimumAddedTime{}--\RQThreeDefaultMaximumAddedTime{}
nanoseconds per call.  The strong wrappers took
\RQThreeStrongMinimumKernelRatio{}--\RQThreeStrongMaximumKernelRatio{} times
as long and added
\RQThreeStrongMinimumAddedTime{}--\RQThreeStrongMaximumAddedTime{}
nanoseconds.  They were faster than the default wrappers for four operations
under CompCert but slower for all six under Clang, showing that their relative
cost depends on the operation and compiler.  The transformation wrappers took
\RQThreeTransformationMinimumKernelRatio{}--\RQThreeTransformationMaximumKernelRatio{}
times as long as their kernels.

Under Clang, the default wrappers took
\RQThreeDefaultMinimumQDRatio{}--\RQThreeDefaultMaximumQDRatio{} times as long
as QD, the strong wrappers took
\RQThreeStrongMinimumQDRatio{}--\RQThreeStrongMaximumQDRatio{} times as long,
and the transformation wrappers took
\RQThreeTransformationMinimumQDRatio{}--\RQThreeTransformationMaximumQDRatio{}
times as long.  These measurements compare total running time and therefore
also include differences between implementations and algorithms.
\FloatBarrier

\section{Related Work}
\label{sec:related-work}

\paragraph{Double-word error analysis.}
The closest numerical basis for our work is the analysis of double-word
addition, multiplication, and division by \citet{Joldes2017} and its Rocq
mechanization by \citet{Rideau2022}.  These results assume that underflow and
overflow do not occur; the mechanization represents this setting using
Flocq's FLX format with an unbounded exponent range.  Earlier results
characterize product corrections near underflow
\citep{BoldoDaumas2003} and establish robust guarantees for \TwoSum{} and
\FastTwoSum{} \citep{Boldo2017}.  We extend these results through the complete
double-word operations under IEEE~754 arithmetic with a finite exponent
range.

\paragraph{Floating-point expansions.}
Floating-point expansions provide additional precision using several values
from an underlying format.  \citet{Shewchuk1997} uses error-free transformations
and nonoverlapping expansions for adaptive geometric predicates, while
\citet{Boldo2017Renorm} formally verify expansion renormalization.  More
recently, \citet{Zhang2025} automatically verify floating-point accumulation
networks and synthesizes an improved network for double-double addition.  Our
work instead considers fixed two-component representations and includes
multiplication and division as well as addition.

\paragraph{Tools for floating-point error analysis.}
Tools in this group derive numerical bounds relating finite-precision
computations to real arithmetic.  Gappa derives bounds and produces proofs
checked in Coq \citep{deDinechin2011}.  FPTaylor uses symbolic Taylor
expansions and global optimization and can emit proof certificates for HOL
Light \citep{Solovyev2015}.  PRECiSA combines static analysis and optimization
with certificates checked in PVS \citep{Titolo2024PRECiSA}.  Daisy provides
modular analysis and optimization for floating-point and fixed-point programs
\citep{Darulova2018}; certificates generated by Daisy can be checked by FloVer
in Coq or HOL4 \citep{Becker2018}.  \citet{Roux2016} formalizes general
techniques for rounding-error proofs, while VCFloat2 performs error analysis
inside Coq \citep{AppelKellison2024}.

\paragraph{Verification of floating-point programs.}
A program proof may use a tool for floating-point error analysis, but must also
connect the resulting numerical analysis to the semantics of the executable
code.  \citet{Ramananandro2016} provide a unified
Coq framework for verifying C programs with floating-point computations.
\citet{AppelBertot2020} organize this connection as layered VST and Flocq
proofs; LAProof and the verified Jacobi solver apply the same structure to
numerical C programs \citep{Kellison2023LAProof,Tekriwal2023Jacobi}.  We follow
this structure, using the IEEE~754 operations and CompCert bridge supplied by
VCFloat2 and adding proofs that the runtime checks establish the conditions
required by the numerical theorems.

\paragraph{Double-double libraries.}
The QD library provides practical double-double and quad-double arithmetic
using algorithms developed by \citet{Hida2001} from techniques introduced by
\citet{Dekker1971} \citep{QD}.  QD does not use runtime checks to establish the
numerical conditions considered in this paper.  We therefore use it as a
comparison implementation in \Cref{sec:evaluation}, while MPFR supplies the
independent reference values.

\section{Conclusion}
\label{sec:conclusion}

\FenceDD{} is a verified C library for double-double arithmetic.
Our Rocq analysis extends earlier double-word results to
IEEE~754 arithmetic, accounting for gradual underflow and
establishing sufficient conditions for finite results.
Using VST, we prove that the C kernels in \FenceDD{} implement the analyzed
algorithms and that successful runtime checks establish the
conditions required by the numerical theorems.

\section*{Data-Availability Statement}
The C sources, Rocq proofs, and evaluation scripts are
ar\-chived \citep{Kellison2026Artifact} and available at
\url{https://github.com/VeriNum/double-double}.



\begin{acks}
The author thanks Andrew Appel and Nate Foster for helpful discussions.
\end{acks}

\bibliography{refs}
\end{document}
